\documentclass[a4paper,12pt,twoside]{memoir}

% Castellano
\usepackage[spanish,es-tabla]{babel}
\usepackage{eurosym}
\selectlanguage{spanish}
\usepackage[utf8]{inputenc}
\usepackage[T1]{fontenc}
\usepackage{lmodern} % scalable font
\usepackage{microtype}
\usepackage{placeins}
\usepackage{dirtree}

\RequirePackage{booktabs}
\RequirePackage[table]{xcolor}
\RequirePackage{xtab}
\RequirePackage{multirow}

% Links
\usepackage[colorlinks]{hyperref}
\hypersetup{
	allcolors = {red}
}

% Ecuaciones
\usepackage{amsmath}

% Rutas de fichero / paquete
\newcommand{\ruta}[1]{{\sffamily #1}}

% Párrafos
\nonzeroparskip
% Huérfanas y viudas
\widowpenalty100000
\clubpenalty100000

% Evitar solapes en el header
\nouppercaseheads



% Imagenes
\usepackage{graphicx}
\newcommand{\imagen}[2]{
	\begin{figure}%[!h]
		\centering
		\includegraphics[width=0.9\textwidth]{#1}
		\caption{#2}\label{fig:#1}
	\end{figure}
	\FloatBarrier
}
\newcommand{\imagenPeque}[2]{
	\begin{figure}%[!h]
		\centering
		\includegraphics[width=0.4\textwidth]{#1}
		\caption{#2}\label{fig:#1}
	\end{figure}
	\FloatBarrier
}
\newcommand{\imagenPequeDos}[2]{
	\begin{figure}%[!h]
		\centering
		\includegraphics[width=0.7\textwidth]{#1}
		\caption{#2}\label{fig:#1}
	\end{figure}
	\FloatBarrier
}
\newcommand{\imagenMasPeque}[2]{
	\begin{figure}%[!h]
		\centering
		\includegraphics[width=0.3\textwidth]{#1}
		\caption{#2}\label{fig:#1}
	\end{figure}
	\FloatBarrier
}

\newcommand{\imagenflotante}[2]{
	\begin{figure}%[!h]
		\centering
		\includegraphics[width=0.9\textwidth]{#1}
		\caption{#2}\label{fig:#1}
	\end{figure}
}



% El comando \figura nos permite insertar figuras comodamente, y utilizando
% siempre el mismo formato. Los parametros son:
% 1 -> Porcentaje del ancho de página que ocupará la figura (de 0 a 1)
% 2 --> Fichero de la imagen
% 3 --> Texto a pie de imagen
% 4 --> Etiqueta (label) para referencias
% 5 --> Opciones que queramos pasarle al \includegraphics
% 6 --> Opciones de posicionamiento a pasarle a \begin{figure}
\newcommand{\figuraConPosicion}[6]{%
  \setlength{\anchoFloat}{#1\textwidth}%
  \addtolength{\anchoFloat}{-4\fboxsep}%
  \setlength{\anchoFigura}{\anchoFloat}%
  \begin{figure}[#6]
    \begin{center}%
      \Ovalbox{%
        \begin{minipage}{\anchoFloat}%
          \begin{center}%
            \includegraphics[width=\anchoFigura,#5]{#2}%
            \caption{#3}%
            \label{#4}%
          \end{center}%
        \end{minipage}
      }%
    \end{center}%
  \end{figure}%
}

%
% Comando para incluir imágenes en formato apaisado (sin marco).
\newcommand{\figuraApaisadaSinMarco}[5]{%
  \begin{figure}%
    \begin{center}%
    \includegraphics[angle=90,height=#1\textheight,#5]{#2}%
    \caption{#3}%
    \label{#4}%
    \end{center}%
  \end{figure}%
}
% Para las tablas
\newcommand{\otoprule}{\midrule [\heavyrulewidth]}
%
% Nuevo comando para tablas pequeñas (menos de una página).
\newcommand{\tablaSmall}[5]{%
 \begin{table}
  \begin{center}
   \rowcolors {2}{gray!35}{}
   \begin{tabular}{#2}
    \toprule
    #4
    \otoprule
    #5
    \bottomrule
   \end{tabular}
   \caption{#1}
   \label{tabla:#3}
  \end{center}
 \end{table}
}

%
%Para el float H de tablaSmallSinColores
\usepackage{float}

%
% Nuevo comando para tablas pequeñas (menos de una página).
\newcommand{\tablaSmallSinColores}[5]{%
 \begin{table}[H]
  \begin{center}
   \begin{tabular}{#2}
    \toprule
    #4
    \otoprule
    #5
    \bottomrule
   \end{tabular}
   \caption{#1}
   \label{tabla:#3}
  \end{center}
 \end{table}
}

\newcommand{\tablaApaisadaSmall}[5]{%
\begin{landscape}
  \begin{table}
   \begin{center}
    \rowcolors {2}{gray!35}{}
    \begin{tabular}{#2}
     \toprule
     #4
     \otoprule
     #5
     \bottomrule
    \end{tabular}
    \caption{#1}
    \label{tabla:#3}
   \end{center}
  \end{table}
\end{landscape}
}

%
% Nuevo comando para tablas grandes con cabecera y filas alternas coloreadas en gris.
\newcommand{\tabla}[6]{%
  \begin{center}
    \tablefirsthead{
      \toprule
      #5
      \otoprule
    }
    \tablehead{
      \multicolumn{#3}{l}{\small\sl continúa desde la página anterior}\\
      \toprule
      #5
      \otoprule
    }
    \tabletail{
      \hline
      \multicolumn{#3}{r}{\small\sl continúa en la página siguiente}\\
    }
    \tablelasttail{
      \hline
    }
    \bottomcaption{#1}
    \rowcolors {2}{gray!35}{}
    \begin{xtabular}{#2}
      #6
      \bottomrule
    \end{xtabular}
    \label{tabla:#4}
  \end{center}
}

% Nuevo comando para tablas grandes con cabecera y filas alternas coloreadas en gris.
\newcommand{\tablaCasosDeUso}[6]{%
	\begin{center}
		\tablefirsthead{
			\toprule
			#5
			\otoprule
		}
		\tablehead{
			\multicolumn{#3}{l}{\small\sl continúa desde la página anterior}\\
			\toprule
			#5
			\otoprule
		}
		\tabletail{
			\hline
			\multicolumn{#3}{r}{\small\sl continúa en la página siguiente}\\
		}
		\tablelasttail{
			\hline
		}
		\bottomcaption{#1}
		\begin{xtabular}{#2}
			#6
			\bottomrule
		\end{xtabular}
		\label{tabla:#4}
	\end{center}
}

%
% Nuevo comando para tablas grandes con cabecera.
\newcommand{\tablaSinColores}[6]{%
  \begin{center}
    \tablefirsthead{
      \toprule
      #5
      \otoprule
    }
    \tablehead{
      \multicolumn{#3}{l}{\small\sl continúa desde la página anterior}\\
      \toprule
      #5
      \otoprule
    }
    \tabletail{
      \hline
      \multicolumn{#3}{r}{\small\sl continúa en la página siguiente}\\
    }
    \tablelasttail{
      \hline
    }
    \bottomcaption{#1}
    \begin{xtabular}{#2}
      #6
      \bottomrule
    \end{xtabular}
    \label{tabla:#4}
  \end{center}
}

%
% Nuevo comando para tablas grandes sin cabecera.
\newcommand{\tablaSinCabecera}[5]{%
  \begin{center}
    \tablefirsthead{
      \toprule
    }
    \tablehead{
      \multicolumn{#3}{l}{\small\sl continúa desde la página anterior}\\
      \hline
    }
    \tabletail{
      \hline
      \multicolumn{#3}{r}{\small\sl continúa en la página siguiente}\\
    }
    \tablelasttail{
      \hline
    }
    \bottomcaption{#1}
  \begin{xtabular}{#2}
    #5
   \bottomrule
  \end{xtabular}
  \label{tabla:#4}
  \end{center}
}



\definecolor{cgoLight}{HTML}{EEEEEE}
\definecolor{cgoExtralight}{HTML}{FFFFFF}

%
% Nuevo comando para tablas grandes sin cabecera.
\newcommand{\tablaSinCabeceraConBandas}[5]{%
  \begin{center}
    \tablefirsthead{
      \toprule
    }
    \tablehead{
      \multicolumn{#3}{l}{\small\sl continúa desde la página anterior}\\
      \hline
    }
    \tabletail{
      \hline
      \multicolumn{#3}{r}{\small\sl continúa en la página siguiente}\\
    }
    \tablelasttail{
      \hline
    }
    \bottomcaption{#1}
    \rowcolors[]{1}{cgoExtralight}{cgoLight}

  \begin{xtabular}{#2}
    #5
   \bottomrule
  \end{xtabular}
  \label{tabla:#4}
  \end{center}
}




\graphicspath{ {./img/} }

% Capítulos
\chapterstyle{bianchi}
\newcommand{\capitulo}[2]{
	\setcounter{chapter}{#1}
	\setcounter{section}{0}
	\chapter*{#2}
	\addcontentsline{toc}{chapter}{#2}
	\markboth{#2}{#2}
}

% Apéndices
\renewcommand{\appendixname}{Apéndice}
\renewcommand*\cftappendixname{\appendixname}

\newcommand{\apendice}[1]{
	%\renewcommand{\thechapter}{A}
	\chapter{#1}
}

\renewcommand*\cftappendixname{\appendixname\ }

% Formato de portada
\makeatletter
\usepackage{xcolor}
\newcommand{\tutor}[1]{\def\@tutor{#1}}
\newcommand{\course}[1]{\def\@course{#1}}
\newcommand{\cotutor}[1]{\def\@cotutor{#1}}
\definecolor{cpardoBox}{HTML}{E6E6FF}
\def\maketitle{
  \null
  \thispagestyle{empty}
  % Cabecera ----------------
\noindent\includegraphics[width=\textwidth]{cabecera}\vspace{1cm}%
  \vfill
  % Título proyecto y escudo informática ----------------
  \colorbox{cpardoBox}{%
    \begin{minipage}{.8\textwidth}
      \vspace{.5cm}\Large
      \begin{center}
      \textbf{TFG del Grado en Ingeniería Informática}\vspace{.6cm}\\
      \textbf{\LARGE\@title{}}
      \end{center}
      \vspace{.2cm}
    \end{minipage}

  }%
  \hfill\begin{minipage}{.20\textwidth}
    \includegraphics[width=\textwidth]{escudoInfor}
  \end{minipage}
  \vfill
  % Datos de alumno, curso y tutores ------------------
  \begin{center}%
  {%
    \noindent\LARGE
    Presentado por \@author{}\\ 
    en Universidad de Burgos --- \@date{}\\
    Tutor: \@tutor{}\\
    	    \@cotutor{}\\
  }%
  \end{center}%
  \null
  \cleardoublepage
  }
\makeatother


% Datos de portada
\title{Herramientas de realidad aumentada orientadas a docencia\\Documentación Técnica}
\author{Saúl Martín Ibáñez}
\tutor{Álvar Arnaiz González}
\cotutor{Pedro Luis Sánchez Ortega}
\date{\today}

\begin{document}

\maketitle



\cleardoublepage



%%%%%%%%%%%%%%%%%%%%%%%%%%%%%%%%%%%%%%%%%%%%%%%%%%%%%%%%%%%%%%%%%%%%%%%%%%%%%%%%%%%%%%%%



\frontmatter


\clearpage

% Indices
\tableofcontents

\clearpage

\listoffigures

\clearpage

\listoftables

\clearpage

\mainmatter

\appendix

\apendice{Plan de Proyecto Software}

\section{Introducción}
 Para la organización del trabajo se ha utilizado GitHub\footnote{\url{https://github.com}} con la extensión de ZenHub\footnote{\url{https://www.zenhub.com/}} para facilitar el seguimiento de las issues gracias a las opciones que ofrece la extensión.
 
En el siguiente enlace se encuentra el repositorio del trabajo: \url{https://github.com/smi0010/TFG_Herramientas_Realidad_Aumentada}.
\section{Planificación temporal}

En un principio, se ha llevado a cabo una reunión a la semana para cada sprint, en la que se exponían los avances realizados durante el sprint de la semana y la planificación del siguiente.

A continuación se explicarán un resumen de los sprints que se han ido desarrollando durante el proyecto.

\subsection{Sprint 0 (15/01/2020-22/01/2020)}
Durante la primera reunión, estuvimos hablando sobre el objetivo del proyecto, de que visión teníamos sobre él, y de hacia donde le podríamos dirigir.
Las tareas que se hicieron fueron sobre la creación del repositorio y la instalación de \LaTeX{} para la documentación. También comencé a investigar sobre diferentes herramientas de realidad aumentada de forma general, para poder compararlas y ver las ventajas y desventajas que tienen.

\subsection{Sprint 1 (22/01/2020-29/01/2020)}
Durante este sprint, hemos acordado ir mejorando y completando  la investigación respecto las herramientas de realidad aumentada que tenía seleccionadas.
Para este sprint decidí centrarme en las herramientas Vuforia y Mergecube.
También como uno de los objetivos del TFG es trabajar con Unity para la creación de ejemplos/proyectos, he incluido una pequeña introducción sobre Unity y los pasos necesarios para su instalación.
\subsection{Sprint 2 (29/01/2020 - 5/02/2020)}
Durante el segundo sprint me centré en documentar sobre las diferentes técnicas de realidad aumentada que se utilizan por la mayoría de herramientas.
También en documentar las herramientas de ArCore y ArKit.

Otro de los puntos que tratamos, fue sobre posibles ideas para la elaboración de un ejemplo en el que poner aprueba las herramientas AR. Surgió la idea de realizar una especie de juego o aplicación <<educativa>> que pueda servir como inicio a la programación. Aunque esta idea nos dimos cuenta que depende de cómo la enfoquemos puede ser demasiado compleja y amplia como para hacer un ejemplo sencillo.
\subsection{Sprint 3 (5/02/2020 - 12/02/2020)}

Me centré en la documentación de las herramientas de Kudan y 8thwall.
Debido que estuve enfermo gran parte de las semana no pude avanzar mas.
\subsection{Sprint 4 (12/02/2020 - 26/02/2020)}
Durante este sript decidí centrarme en las herramientas de Wikitude, OpenCV, ZapWorks, y en realizar también correcciones y mejoras de la documentación que el tutor había señalado.
También estuve pensando en posibles ideas para realizar un ejemplo aplicando las herramientas de realidad aumentada. 
\subsection{Sprint 5 (26/02/2020 - 04/03/2020)}
Se estuvo hablando respecto al ejemplo, para empezar ya con el. Sobre la idea de hacer un pequeño juego en el que hay un <<laberinto>> y que el usuario tuviera que resolver el camino moviéndose por comandos de una forma similar a una programación por blocking pero simplificado, todo esto aplicado con las herramientas de realidad aumentada. 

Dado que aún la idea parecía algo incompleta y complicada de ajustar a las herramientas de realidad aumentada, vamos a intentar primero un ejemplo mas sencillo, este sería uno en la que poder ver las fases de crecimiento en una <<planta>>. 
Durante el sprint probare a poder usar la cámara de forma que sea como una regadera, así cuando este encima del Mercube donde estaría la planta virtual, se simule en realidad aumentada, que cae agua desde la cámara del teléfono hacia la planta. Otra forma pensada mas simplificada, sería simplemente añadir en pantalla un botón, que al presionarle se simule que cae agua sobre la planta sin tener que mover el teléfono.

\subsection{Sprint 6 (04/03/2020 - 11/03/2020)}

Durante la reunión de este sprint mostré los primeros pasos que había realizado del ejemplo a desarrollar en Unity para aplicar herramientas AR. También comenté mis ideas sobre como continuar y los puntos en que me quería centrar durante este sprint.

Centrarme en poder detectar las partículas que simulan ser agua, cuando estas colisionan con un objeto determinado. Y que contando el número de contactos poder determinar la progresión para poder pasar a otra etapa o nivel. También hacer uso del sensor giroscopio del móvil, para determinar con la inclinación del dispositivo la cantidad de partículas que caen.

Como la idea del ejemplo sería seguir el crecimiento de una planta en realidad aumentada desde un punto educativo, pensamos en que aparte de poder echarle agua, también se podría cambiar entre diferentes opciones, como abonos, sulfatos, etc. También se podrían considerar establecer normas del tipo, escoger la temporada en la que se planta, si necesita que le quiten malas hiervas, etc. Por lo que habría que incluir una forma de cambiar entre esas diferentes opciones.

Otro de los aspectos que he estado mirando, es mejorar el posicionamiento virtual del objeto 3D en el marcador. Ya que con la primera configuración que he probado, cuando se mueve la cámara (el móvil) el objeto se mueve de su posición demasiado y no queda estable.

\subsection{Sprint 7 (11/03/2020 - 18/03/2020)}
Durante la reunión de este sprint mostré los avances que había realizado sobre el ejemplo de la aplicación. También les expliqué que estuve realizando varias pruebas añadiendo las funcionalidades del SDK de MergeCube en la aplicación. Consiguiendo mejorar los resultados con el uso conjunto de Vuforia + Mergecube, que únicamente con Vuforia.

Otro de los aspectos tratados fue sobre la herramienta de CoSpace, la necesidad de explicar las ventajas de su uso, para explicar su viabilidad económica.

Durante la semana realicé algunos avances en la documentación, y estuve solucionando un bug en el uso del giroscopio. El bug en cuestión se traba de que el eje de gravedad de las partículas, no coincidía con el eje de gravedad normal. El problema resultó ser que, el parámetro <<Simulation Space>> debía de estar con la opción Word. Otro problema de este bug, era que el eje Z, no tenía la interpretación deseada para nuestro ejemplo. Para esto en el script del giroscopio, se modificó para invertir el eje Z.
 
\subsection{Sprint 8 (18/03/2020 - 25/03/2020)}
En este Sprint, la reunión se tuvo que realizar por videoconferencia, debido a la cuarentena en la que se encontraba el país durante esos días. Expliqué los avances que había realizado en la documentación y algunas de las dudas que me habían surgido sobre como debería avanzar en ciertos puntos de la documentación. También expliqué sobre el bug que había corregido en la aplicación, aunque no puede mostrárselo ya que al ser la reunión online  no pude preparar un vídeo para enseñarlo.
 
Para este sprint seguí avanzando en la documentación, me centré en especificar los requisitos y objetivos de la aplicación. También diseñé en unas imágenes un prototipo de la interfaz. En la aplicación trabajé en añadir diferentes fases del crecimiento de un planta, que cuando se completa el progreso de una fase salta a la siguiente. 
\subsection{Sprint 9 (25/03/2020 - 01/04/2020)}
Durante el sprint de esta semana me he centrado en dotar a la aplicación de un prototipo de la interfaz y de menú funcional. En la interfaz el principal punto en que me he centrado, ha sido el de realizar una barra de progreso, la cual se rellena a medida que el nivel es completado. En el menú, me he centrado entre la navegación entre los diferentes niveles del menú, y el más importante, que desde el menú se pueda inicializar el juego correctamente.
También se ha investigado sobre la posibilidad de que la aplicación mande un fichero o un mensaje, a un servidor o pagina web. Esto serviría para realizar un encuesta de opinión de la aplicación sin tener que salir de ella. 
\subsection{Sprint 10 (01/04/2020 - 15/04/2020)}
Como objetivo durante el sprint tenía mejorar la interfaz, principalmente con la opción de poder activar y desactivar las herramientas del juego y añadir ventanas emergentes en las que se podría leer información de ayuda o complementaria para el usuario. También centrarme en mejorar algunos de los apartados en la documentación.
Otro de los puntos era el hacer pruebas en CoSpaces, intentando realizar un ejemplo similar al que se está desarrollando en Unity y poder hacer comparaciones.


\section{Estudio de viabilidad}

\subsection{Viabilidad económica}

En este apartado se evaluara la viabilidad económica de las herramientas que se han investigado.


CoSpaces, ofrece varias tarifas aparte de la básica que es gratuita, la más barata de 69,99\euro{}  anuales y la mas cara es de 1349,99\euro{} anuales.

La tarifa básica, ofrece la creación de una clase con capacidad para 30 personas (alumnos y profesores). En esta clase virtual, el profesor puede crear unicamente una tarea y asignarse al alumnado de forma individual o por grupos. El número máximo de profesores por clase en la tarifa básica es de uno.

La tarifa básica, permite la creación de un máximo de dos espacios, es decir un máximo de dos proyectos. En la creación de ejemplos, estos se encuentran limitados a unos pocos modelos 3D de CoSpace, y a un limitado número de opciones de programación por CoBlocks. En la tarifa básica no es posible realizar proyectos para el MergeCube.

Con la tarifa Pro de 69,99\euro{}, el número de personas máximas para una clase, aumenta en 5, permitiendo también que haya mas de un profesor en la clase. Permite crear mas de una clase. En el apartado de espacios, la tarifa Pro, permite crear mas de dos espacios, y en estos espacios da acceso a todos los modelos 3D de CoSpace, y permite utilizar en la programación todos los bloques de CoBlocks y scripts.
Si se desea utilizar el Meregecube, la licencia asciende hasta los 79,99\euro{}. 

El resto de tarifas que posee CoSpaces, su principal característica es el número de personas simultaneas para una clase que puede tener, subiendo en la última tarifa hasta 400 personas adicionales.

La tarifa que sería recomendable se trata de la primera ofertada en el plan Pro, ya que aunque se vaya a utilizar con un número de alumnos superior al ofertado, sería posible organizarse por turnos.

\subsubsection{Costes Personal}
El proyecto ha sido desarrollado por una única persona, durante un periodo de seis meses a tiempo parcial.

\subsubsection{Costes Material}
Mergecube...
\subsubsection{Costes Totales}
\subsection{Viabilidad legal}



\apendice{Especificación de Requisitos}

\section{Introducción}
En el presente anexo se detallan los objetivos generales de la aplicación desarrollada y el catálogo de requisitos funcionales y no funcionales así como los diagramas de casos de uso.
\section{Objetivos generales}
\begin{itemize}
	\item Un objetivo es el de estudiar el estado del arte en este tipo de aplicaciones de realidad aumentada. 
	\item Realizar una aplicación para el ámbito educativo, empleando las aplicaciones AR estudiadas.
	\item Utilizar MergeCube para el diseño e implementación de la aplicación.
	\item Como ejercicio se ha diseñado el crecimiento de una planta, en el que se representaría las principales fases del crecimiento. Sería para usar en clases de biología de primaria/secundaria.
	\item Que pueda usarse en el mayor número de dispositivos compatible posibles.
\end{itemize}
\section{Catálogo de requisitos}
\subsection{Requisitos Funcionales}
 \begin{itemize}
	\item \textbf{RF 1:} Crear una aplicación educativa.
	\begin{itemize}
	\item \textbf{RF 1.1:} Permitir que el usuario configure el nivel.
	\item \textbf{RF 1.2:} 
	\item \textbf{RF 1.3:} Ofrecer ayuda o tutorial.	
	\item \textbf{RF 1.4:} Realizar capturas de pantalla y guardarlas en el dispositivo.
	\item \textbf{RF 1.5:} Permitir realizar encuestas de valoración.
	\item \textbf{RF 1.6:} Acceder a la información de los sensores del dispositivo.
	
	\end{itemize} 
	\item \textbf{RF 2:} Modelar un ejercicio de realidad aumentada sobre el crecimiento de una planta.
	\begin{itemize}
	\item \textbf{RF 2.1:} Representar visualmente el avance en cada fase.
	\item \textbf{RF 2.2:} Representar visualmente las herramientas con las que se avanza en el juego, simulando que el dispositivo puede manejar esa herramienta.
	\item \textbf{RF 2.3:} Cada fase de un nivel, estará representado por una etapa del crecimiento de una planta.
	
	\end{itemize} 
	
	
\end{itemize}
\subsection{Requisitos no Funcionales}
\begin{itemize}
	\item \textbf{RNF 1:} La aplicación debe de resultar fácil de usar y intuitiva para el usuario.
	\item \textbf{RNF 2:} Debe tener un rendimiento aceptable, para un uso cómodo.
	\item \textbf{RNF 3:} Debe de ser responsiva, pudiéndose adaptar la interfaz, a la pantalla de cualquier dispositivo.
\end{itemize}
\section{Especificación de requisitos}

\imagen{Anexos/casosdeuso}{Diagrama de casos de uso.}
\subsection{Descripción de los casos de uso}
%\tablaCasosDeUso{Caso de Uso 1: Menu}{}{}{}
\subsubsection{Caso de uso 1: Configurar nivel.}
\begin{itemize}
	\item \textbf{Descripción:} Configurar los aspectos básicos del juego en la aplicación.
	\item \textbf{Precondiciones:} Encontrarse en el menú principal.
	\item \textbf{Requisitos:} RF-1 RF1.1
	\item \textbf{Acciones:}
	\begin{enumerate}
		\item Seleccionar en el menú el botón Opciones.
		\item Cambiar a la pantalla del menú de opciones.		
		\item Escoger una de las opciones a configurar.		
	\end{enumerate}
	\item \textbf{PostCondiciones:} Se ha entrado en la configuración correctamente. al salir los cambios se han guardado y hecho efectivos.
	\item \textbf{Excepciones:} -
\end{itemize}
\subsubsection{Caso de uso 1.1: Escoger modo target.}
\begin{itemize}
	\item \textbf{Descripción:} Escoger el tipo de target que se usará en la aplicación.
	\item \textbf{Precondiciones:} haber entrado en opciones
	\item \textbf{Requisitos:} RF-1 RF-1.1
	\item \textbf{Acciones:}
	\begin{enumerate}
		\item Elegir una de las opciones disponibles.
		\item Configurar parámetros a decisión del usuario.
	\end{enumerate}
	\item \textbf{PostCondiciones:} guardar los cambios realizados.
	\item \textbf{Excepciones:} -
\end{itemize}
\subsubsection{Caso de uso 2:  Seleccionar modo juego AR}
 \begin{itemize}
 	\item \textbf{Descripción:} Seleccionar uno de los modos de juego disponibles.
 	\item \textbf{Precondiciones:} Encontrarnos en el menú principal.
 	\item \textbf{Requisitos:} RF-1
 	\item \textbf{Acciones:}
 	\begin{enumerate}
 		\item Seleccionar en el menú el botón Modos de Juego.
 		\item cambiar a la pantalla del menú de modos de juego.
 	\end{enumerate}
 	\item \textbf{PostCondiciones:} Acceder al menú de modos de juego.
 	\item \textbf{Excepciones:}	-
 \end{itemize}
\subsubsection{Caso de uso 2.1: Seleccionar semilla.}
\begin{itemize}
	\item \textbf{Descripción:} El usuario escogerá una de las semilla/planta disponibles y se iniciara el modo AR de esa semilla.
	\item \textbf{Precondiciones:} Haber seleccionado el modo de las plantas.
	\item \textbf{Requisitos:} RF-1 RF-1.6 RF-2
	\item \textbf{Acciones:}
	\begin{enumerate}
		\item Seleccionar en una de las disponibles.
		\item Iniciar el modo AR sobre la semilla seleccionada.
	\end{enumerate}
	\item \textbf{PostCondiciones:} Se ha iniciado el modo de realidad aumentada correctamente.
	\item \textbf{Excepciones:} Si se trata de iniciar una planta aun no implementada, no se iniciará el modo AR.
\end{itemize}
\subsubsection{Caso de uso 2.1.1: Escoger una herramienta}
\begin{itemize}
	\item \textbf{Descripción:} Escoger entre una de las diferentes herramientas <<jugables>> del modo AR.
	\item \textbf{Precondiciones:}
	\begin{enumerate}
		\item Se ha iniciado el modo AR.
		\item Sensor de la cámara encendido.
	\end{enumerate}
	\item \textbf{Requisitos:} RF-2 RF-2.2
	\item \textbf{Acciones:} 
	\begin{enumerate}
		\item Pulsar  botón de la esquina inferior derecha.
		\item Se desplegará una barra con botones de distintas herramientas.
		\item Escoger una de las herramientas y pulsar.
		\item Usar la harramienta seleccionada y completar sus funciones en el nivel.
	\end{enumerate}
	\item \textbf{PostCondiciones:} La herramienta sale en pantalla y es funcional.
	\item \textbf{Excepciones:} 
	\begin{enumerate}
		\item La funcionalidad de la herramienta ya ha sido completada.
		\item El nivel ya ha sido completado.
	\end{enumerate}

\end{itemize}
\subsubsection{Caso de uso 2.1.2: Captura de pantalla. }
 \begin{itemize}
 	\item \textbf{Descripción:} Realiza una captura de pantalla de ese instante y la guarda en la galería del terminal.
 	\item \textbf{Precondiciones:}
	 		\begin{enumerate}
	 		\item Se ha iniciado el modo AR
	 		\item Sensor de la cámara encendido.
	 		
	 	\end{enumerate}
 	\item \textbf{Requisitos:} RF-1.4 RF-2 RF-2.1
 	\item \textbf{Acciones:}
	 	\begin{enumerate}
	 		\item Seleccionar captura de pantalla.
	 		\item Seleccionar captura de pantalla.
	 		\item Guardar en un archivo png.
	 		\item Volver al funcionamiento
	 	\end{enumerate}
 	
 	\item \textbf{PostCondiciones:}
	 	\begin{enumerate}
	 		\item Captura guardada correctamente como un archivo de imagen.
	 		\item El modo AR sigue funcionando y no se interrumpe.
	 		
	 	\end{enumerate}
 
 	\item \textbf{Excepciones:}  No hay espacio suficiente de almacenamiento.
 \end{itemize}
\subsubsection{Caso de uso  2.1.3: Visualizar interfaz.}
\begin{itemize}
	\item \textbf{Descripción:} La interfaz del modo AR se hace visible para el usuario.
	\item \textbf{Precondiciones:} Entrar en el modo jugable AR
	\item \textbf{Requisitos:} RF-2 RF-2.1 RF-2.2 RF-2.3
	\item \textbf{Acciones:}
		\begin{enumerate}
			\item Entrar en el modo AR.
			\item Cargar interfaz
		\end{enumerate}
	\item \textbf{PostCondiciones:}
		\begin{enumerate}
			\item La interfaz ha sido cargada correctamente.
		\end{enumerate}
	\item \textbf{Excepciones:} Error al cargar el modo AR.
\end{itemize}
\subsubsection{Caso de uso 3: Ayuda.}
\begin{itemize}
	\item \textbf{Descripción:} Proporciona una guía o ayuda básica a los usuarios, respecto al uso de la aplicación.
	\item \textbf{Precondiciones:} -
	\item \textbf{Requisitos:} RF-1 RF-1.3
	\item \textbf{Acciones:}
	\begin{enumerate}
		\item Seleccionar en el menú el botón Ayuda.
		\item Cambiar a la pantalla del menú de Ayuda.
	\end{enumerate}
	\item \textbf{PostCondiciones:} La ventana que muestra la información de ayuda, ha cargado correctamente.
	\item \textbf{Excepciones:} -
\end{itemize}
\subsubsection{Caso de uso 4: Encuestas}
\begin{itemize}
	\item \textbf{Descripción:}  Al completarse un nivel se da la posibilidad de realizar una encuesta sobre la app.
	\item \textbf{Precondiciones:} Completar un nivel.
	\item \textbf{Requisitos:} RF-1 RF-1.5
	\item \textbf{Acciones:}
	\begin{enumerate}
		\item Finalizar un nivel.
		\item Entrar en link de la encuesta.
		\item Redirigir a la página, usando el navegador del dispositivo.
		\item Contestar las preguntas.
	\end{enumerate}
	\item \textbf{PostCondiciones:} Se ha redirigido a la pagina de la encuesta..
	\item \textbf{Excepciones:} No hay conexión a Internet, por lo que no se puede acceder a la encuesta.
	
\end{itemize}

\subsection{Actores}
\begin{itemize}
	\item \textbf{Usuario:} Usuario que hará uso aplicación.
\end{itemize}



\apendice{Especificación de diseño}

\section{Introducción}
 Para este apartado se presentarán los diseños que se han seguido para el desarrollo de la aplicación.
 

\section{Diseño procedimental}
En la figura  \ref{fig:Diagrama_secuencia} se muestra un diagrama de secuencia, en el que se puede observar el comportamiento básico de la aplicación a la hora de buscar, detectar un marcador y renderizar su modelo asociado. Para crear el diagrama se ha usado la herramienta online Draw.io\footnote{\url{https://app.diagrams.net/}}.


\imagen{Anexos/Diagrama_secuencia}{Diagrama de secuencia del proceso de buscar un marcador.}

\section{Diseño arquitectónico}


Dado que se trata de una aplicación desarrollado en Unity la estructura que sigue viene dada por defecto propiamente por Unity. En la imagen\ref{fig:Anexos/unity_xr_new_architecture} podemos ver una sencilla representación de la estructura de Unity con respecto al plugin de Vuforia. 

La plataforma XR de Unity se trata de una capa dedicada específicamente al control de tecnologías de realidad aumentada y realidad virtual, con el objetivo conseguir una mayor compatibilidad y facilidades de desarrollo para software de extensión. Hasta la versión 2019.3 de Unity, Vuforia contaba con soporte integrado en Unity por defecto, apartir de dicha versión es necesario descargar Vuforia Engine desde el Portal de desabolladores de Vuforia. 

 
El sdk de Vuforia interactua directamente con Unity XR, para el uso de los recursos de realidad aumentada (Cámara, targets...). 

El Unity Core será el encargado del renderizado general de los recursos y demás funciones de Unity.

Por ultimo el sdk de Android interactuara con Unity para realizar la compilación del APK y poder un correcto porte para una aplicación Android.

\imagen{Anexos/unity_xr_new_architecture}{Arquitectura XR de Unity.}




\section{Diseño de interfaz}

En las siguientes imágenes se pueden ver, el prototipo de interfaz pensado para la aplicación. 

En las figuras \ref{fig:Anexos/Interfaz/page1}, \ref{fig:Anexos/Interfaz/page2} y \ref{fig:Anexos/Interfaz/page2_2} se muestra un prototipo de los menús de la aplicación. En la segunda, se puede escoger diferentes modos de juego, por el momento en esta aplicación unicamente se está desarrollando un modo de juego. 
Otra posibilidad para los menús, sería integrarlos por realidad aumenta en el MergeCube, teniendo que girarle y jugar con él para navegar por el menú.

En las figuras \ref{fig:Anexos/Interfaz/page3} y \ref{fig:Anexos/Interfaz/page4}, se muestra un prototipo de la interfaz del juego. Tendríamos botones para seleccionar las diferentes herramientas posibles en la parte inferior. Y en la parte superior una barra de progreso, para indicar el progreso de cada fase. Para regresar al menú principal se podría usar los botones del terminal o añadir uno a la interfaz también.

\imagenPeque{Anexos/Interfaz/page1}{Prototipo del menú principal de la aplicación.}

\imagenPeque{Anexos/Interfaz/page2}{Prototipo del menú de juegos de la aplicación.}

\imagenPeque{Anexos/Interfaz/page2_2}{Prototipo del menú de tipos de plantas disponibles.}

\imagenPeque{Anexos/Interfaz/page3}{Prototipo de la interfaz en modo AR.}
A la hora de probar en distintos dispositivos, en los dispositivos que tienen en la pantalla notch(cámara frontal) se aprecio que la interfaz no tenia en cuenta el notch, por lo que este tapaba una parte de la barra de progreso. Para corregirlo se ha bajado la barra de progreso.

\imagenPeque{Anexos/Interfaz/page4}{Prototipo de la interfaz en modo AR.}
\apendice{Documentación técnica de programación}

\section{Introducción}

\section{Estructura de directorios}

\section{Manual del programador}
\subsection{Unity}
Instalación.

Para Instalar el motor Unity es necesario descargar Unity Hub: \url{https://store.unity.com/es/download-nuo}

Unity Hub se trata de un launcher desde el que se pueden descargar varias versiones de Unity simultáneamente, pudiendo escoger la que mejor se acomode a las necesidades del usuario. También ofrece una serie de tutoriales para iniciarse en el desarrollo de Unity. Por último ofrece un listado de los proyectos del usuario, así como a que versión pertenecen, pudiendo escoger con que versión de Unity desean ejecutarlos.

En la pestaña de installs, pulsando el botón <<Add>>, saldrá una ventana donde poder escoger la versión deseada, al pasar el siguiente paso deberemos escoger los módulos complementarios para Unity, por el momento sería necesario el de Android, para poder pasar nuestros proyectos a una aplicación Android, importante desplegar las opciones del módulo y seleccionar ambas, pues para la compilación de una aplicación es estrictamente necesario el APK de Android.\imagen{unity01}{Paso 1: Seleccionar la versión de Unity }
\imagen{unity02}{Paso 2: Seleccionar}

Una vez terminada la instalación ya es posible comenzar a crear proyectos. 
Para poder descargar y usar Assets desde la tienda de Unity es necesario tener una cuenta de usuario.

\section{Compilación, instalación y ejecución del proyecto}

\section{Pruebas del sistema}

\apendice{Documentación de usuario}

\section{Introducción}
 En este apartado se explicarán, los requisitos que necesita el usuario para ejecutar la aplicación, los pasos para su instalación, y cómo usarla.
  
\section{Requisitos de usuarios}
 El ejemplo desarrollado se trata de una aplicación Android, por lo que el primer requisito necesario es disponer de un Dispositivo Android, con una versión Android 6.0 (Marshmallow) o superior.
  
 Ademas deberá cumplir los siguientes requisitos\cite{vuforia_devices}:
 \begin{itemize}
 	\item CPU Quad Core.
 	\item 1-2GB de memoria RAM.
 	\item Sensor de Giróscopio.
 	\item Se recomienda una cámara que permita grabar mínimo en 720p.  
 \end{itemize}
También es necesario disponer de un MergeCube. Para utilizar el Mergecube tenemos las opciones de comprar el original o construirse uno de papel. En ambos casos se han probado y son validados, la pega del de papel es su fragilidad y que dependiendo de la calidad de la impresión los resultados pueden variar.

\section{Instalación}

Dado que la aplicación no se ha publicado en Play Store, se deberá descargar el APK desde (\href{https://github.com/smi0010/TFG_Herramientas_Realidad_Aumentada/releases/tag/UBU_AR}{UBU AR}).
La instalación de la aplicación es sencilla, como la de cualquier otra APK. Una vez descargado el APK en el dispositivo, ejecutarle y seguir los pasos. 

Dado que la aplicación no se ha descargado desde una fuente oficial, como podría ser la Play Store, es posible que en el primer intento de instalación, el sistema solicite habilitar la opción de instalar aplicaciones de origen desconocido, que por defecto en los dispositivos viene desactivada.
En un principio los propios pasos de instalación permite habilitar la opción. En caso contrario deberás buscar en opciones del sistema. En caso de tener complicaciones para encontrar la ubicación de dicha opción, es recomendable realizar una sencilla búsqueda en Google, de ``Activar la instalación de aplicaciones de fuentes desconocidas en <<Modelo del dispositivo usado>>, con <<Versión de Android que posee>>``.


\section{Manual del usuario}

En este apartado se explicará los pasos básicos para aprender a usar la aplicación.

En la pantalla principal, al ejecutar la aplicación, encontraremos un menú con 4 cuatro opciones: 
\begin{itemize}
	\item Modos de Juego
	\item Opciones
	\item Ayuda
	\item Info
\end{itemize}


En la primera opción ``\textbf{Modos de Juego}'' el usuario podrá buscar entre los diferentes niveles disponibles la aplicación y seleccionar uno, ejecutando así el modo jugable. en la imagen (ref).

En ``\textbf{Opciones}'', se puede configurar diferentes aspectos de la aplicación, como el tipo de target que se usará en la realidad aumentada. En la imagen~(\ref{fig:Anexos/appOpciones}) podemos ver las opciones disponible.

\imagen{Anexos/appOpciones}{Imagen perteneciente a las opciones de la aplicación.}

En la opción ``\textbf{Ayuda}'', se ofrece una guía/tutorial en la que se explican los pasos básicos para el uso de la aplicación y poder resolver las dudas de los usuarios que se puedan encontrar.

La cuarta y ultima Opción se trata de una venta en la que se muestra un resumen respecto al TFG y la aplicación. Se detallan datos como el autor, lo tutores del TFG. 


\subsection{Modo MergeCube }

Al seleccionar uno de los modelos de planta teniendo seleccionado el modo con el MergeCube (ver figura~\ref{fig:Anexos/appOpciones}), la aplicación cambiará al modo de realidad aumentada activando la cámara del dispositivo. Ahora tendremos una interfaz que mostrará una barra de progreso en la parte superior, un contador del tiempo restante para acabar la fase en la parte superior derecha. En la parte inferior derecha hay un botón, que al pulsarse, despliega una barra con las herramientas disponibles para completar las fases. En la imagen \ref{fig:Anexos/appInterfaz} podemos ver unas muestras de la interfaz.

En este modo, una vez que la cámara ya está activa, al enfocar el Mergecube, será detectado y se cargará el modelo de realidad aumentada. En este punto es importante recordar, que para una correcta detección de un marcador, se precisa de estar en un ambiente con una iluminación adecuada para que el marcador no se vea afectado por sombras o focos de luz que reduzcan su visibilidad para el sensor de la cámara.

\imagen{Anexos/appInterfaz}{Imagen perteneciente a la aplicación, con la interfaz desplegada.}

\subsection{Modo CapturaTarget}

En el caso de tener seleccionado el \textbf{modo de capturar target}~(ver figura~\ref{fig:Anexos/appOpciones}), la interfaz contará con dos elementos adicionales: 
\begin{itemize}
	\item El primero se trata de una barra de colores (rojo, amarillo, verde) que indica la calidad óptima para que el objetivo enfocado sea un target. En la imagen~\ref{fig:Anexos/appCapturaTarget} se puede observar como al apuntar a una fondo blanco indica que tiene mala calidad, y al apuntar un fondo con una composición de elementos y colores buena la indica que es optimo para ser un target.
	\item El segundo elemento será el disparador situado en la parte inferior central. Una vez que tengamos un objetivo que cumpla los requisitos de calidad, podremos presionar el disparador, de manera que el target se guardará temporalmente (al cerrar la aplicación o volver al menú los target se reinician). 
\end{itemize}

Cuando capturamos un target, el disparador y la barra de colores desaparecerán y saltará un mensaje preguntándonos si deseamos continuar con dicho target. Si seleccionamos que ``no``, tendremos que volver a capturar un nuevo target. Sobre este target capturado aparecerá el modelo 3D correspondiente y ahora se podrá ``jugar`` del mismo modo que se hacia con el Mergecube. 

En el modo jugable el número de target simultáneos esta limitado a uno, si se desea capturar un nuevo target se deberá indicar presionando el botón de la parte inferior izquierda, donde nos preguntará si deseamos continuar con el target, al presionar que ``no``, volverá a salir las interfaz de captura para tomar un nuevo target. Al capturar un nuevo target sustituirá al anterior, pero podremos seguir con el mismo progreso del nivel que teníamos con el anterior target. 

\imagen{Anexos/appCapturaTarget}{Imagen pertenecientes a la captura de un target.}

\subsection{Barra de Herrameintas}
En la barra de herramientas contamos con dos elementos, una regadera y un paquete de abono. La primera es un icono de gotas de agua, al seleccionar nos saldrá en pantalla una regadera que automáticamente echará agua, el objetivo es mover el dispositivo hasta conseguir que el agua caiga en el Mergecube. Cuando el agua colisione con el modelo proyectado en el Mergecube, la barra de progreso subirá hasta completarse. 

El segundo elemento, se trata de un bote de fertilizante, al pulsarse saldrá en la pantalla un botón que al ser presionado saldrá delante como si se hubiese lanzado contra el Mergecube.
Para completar una fase es necesario haber completado ambos casos anteriores descritos.

También en la barra de herramientas aparecen dos opciones, la de Activar el Device Tracker con el botón ``Fijar posición``, y por último la opción de sacar una captura de pantalla.

\imagenPeque{Anexos/barraHerramienta}{Imagen perteneciente a la barra de herramientas disponibles.}

Actualmente los modos de juego de las plantas constan de cinco fases~(ver imagen~\ref{fig:Anexos/appFases}), cada una correspondiente con un estado de la planta (semilla, brote, planta pequeña, planta mediana, planta final). Para completar una fase, es necesario, completar la barra de progreso de agua, y haber echado fertilizante. Si el fertilizante se ha utilizado en una fase se indica con una luz verde en la barra de progreso.

\imagen{Anexos/appFases}{Imagen perteneciente a las fases de la planta de tomate.}

\subsection{Modo Múltiples Marcadores}
En el modo de captura target, también cuenta con un modo de múltiples marcadores. En este modo no se puede ``jugar``, pero es posible tener múltiples marcadores personalizados al mismo tiempo. Esto da la posibilidad de poder observar varios los modelos de la aplicación al mismo tiempo, como podemos ver en la imagen~\ref{fig:Anexos/appMultiTarget}. Por defecto la aplicación esta limitada a cinco marcadores simultáneos como máximo, para asegurar el rendimiento. En caso de realizar una sexta captura borrará la primera.

\imagenMasPeque{Anexos/appMultiTarget}{Imagen perteneciente al modo de múltiples marcadores simultáneos.}

En este modo podemos cambiar el tamaño de los modelos 3D, para ello primero debemos tocar el modelo que queramos modificar. Y para modificar unicamente hay que pulsar los botones de la parte interior izquierda.
En las opciones de la parte inferior derecha, podemos seleccionar que modelo 3D queremos que se asocie al marcador que vamos a capturar.


 




\bibliographystyle{plain}
\bibliography{bibliografiaAnexos}

\end{document}
